%%%%%%%%%%%%%%%%%%%%%%%%%%%%%%%%%%%%%%%%%
% The Legrand Orange Book
% LaTeX Template
% Version 2.3 (8/8/17)
%
% This template has been downloaded from:
% http://www.LaTeXTemplates.com
%
% Original author:
% Mathias Legrand (legrand.mathias@gmail.com) with modifications by:
% Vel (vel@latextemplates.com)
%
% License:
% CC BY-NC-SA 3.0 (http://creativecommons.org/licenses/by-nc-sa/3.0/)
%
% Compiling this template:
% This template uses biber for its bibliography and makeindex for its index.
% When you first open the template, compile it from the command line with the 
% commands below to make sure your LaTeX distribution is configured correctly:
%
% 1) pdflatex main
% 2) makeindex main.idx -s StyleInd.ist
% 3) biber main
% 4) pdflatex main x 2
%
% After this, when you wish to update the bibliography/index use the appropriate
% command above and make sure to compile with pdflatex several times 
% afterwards to propagate your changes to the document.
%
% This template also uses a number of packages which may need to be
% updated to the newest versions for the template to compile. It is strongly
% recommended you update your LaTeX distribution if you have any
% compilation errors.
%
% Important note:
% Chapter heading images should have a 2:1 width:height ratio,
% e.g. 920px width and 460px height.
%
%%%%%%%%%%%%%%%%%%%%%%%%%%%%%%%%%%%%%%%%%

%----------------------------------------------------------------------------------------
%    PACKAGES AND OTHER DOCUMENT CONFIGURATIONS
%----------------------------------------------------------------------------------------

\documentclass[11pt,fleqn,table]{book} % Default font size and left-justified equations

%----------------------------------------------------------------------------------------

\input{./structure_temp.tex} % Insert the commands.tex file which contains the majority of the structure behind the template

\begin{document}


%-----------------------------------------------------------------------------------------------------------------------------------------
%
% Create Amazing Worksheets and Exams (Video 7: Putting a margin box around every page)
% 
% Function:
%     
%    https://www.youtube.com/watch?v=qVZGhgatyRk&list=PLg5nrpKdkk2DWcg3scb75AknF7DJXs8lk&index=8
%
% Date: 2025.08.13
%
%---------------------------------------------------------------------------

% \begin{tcolorbox}[%
    % enhanced,
	% breakable,
	% frame hidden,
	% overlay broken={
	    % \draw[line width=0.5mm,
		      % gray,
			  % rounded corners]
			      % (frame.north west) rectangle (frame.south east);},
	% colback=white,
    % height fixed for=all]{}
	
\pagestyle{pgstyPaper} % Print headers again
\fancyhf{}

\input{./Var/Math_9709_ms.tex}

% \end{tcolorbox}




\end{document}
